\documentclass[11pt]{article}
\usepackage[utf8]{inputenc}
\usepackage[margin=1in]{geometry}
\usepackage{hyperref}
\usepackage{booktabs}
\usepackage{color}
\usepackage{amsmath}
\usepackage{listings}

\definecolor{primary}{RGB}{30, 41, 59}
\definecolor{secondary}{RGB}{79, 70, 229}
\definecolor{teal}{RGB}{13, 148, 136}

\lstset{
  basicstyle=\ttfamily\small,
  keywordstyle=\color{secondary}\bfseries,
  commentstyle=\color{gray},
  stringstyle=\color{teal},
  breaklines=true
}

\title{\textbf{Pipeline-Based Speaker Guide: GeoPhoto}\\
\large Unified Slide Walkthrough (Slides 1-39) \& Explanations for HyLight}
\author{Prepared for: Rares Olteanu}
\date{June 2026}

\begin{document}

\maketitle
\newpage
\tableofcontents
\newpage

\section{Introduction: The Pipeline Narrative}
To present like a professional engineer, you must tell a **story**. 
Instead of talking about code in isolation, you will walk Jules through the \textbf{Lifecycle of an Aerial Image} captured by HyLight's airship, showing how it flows from the sensor, through your frontend and backend, gets processed by AI, goes through test suites, and scales to the cloud.

This guide gives you the exact script and explanation for every slide (1 to 39) in this pipeline.

---

\section{Pipeline Walkthrough & Scripts}

\subsection{Slide 1: Title Slide (The Hook)}
\begin{itemize}
    \item \textbf{Focus:} Introduce yourself and your pipeline topic.
    \item \textbf{Speaker Script:} 
    \textit{``Hi everyone, I'm Rares Olteanu. Today, I'm presenting the Full-Stack Media Ingestion Pipeline. We will walk through the lifecycle of a geospatial image captured by an inspection sensor—explaining how it gets processed, secured, and validated in our codebase, and how it scales to production levels.''}
\end{itemize}

\subsection{Slide 2: The Pipeline Flow (The Agenda)}
\begin{itemize}
    \item \textbf{Focus:} Set the expectations.
    \item \textbf{Speaker Script:} 
    \textit{``We will trace the image through 6 distinct pipeline stages: client-side sensor ingestion, backend gateway validation, background AI enrichment, security guardrails, automated test suites, and finally, database scaling. This maps directly to industrial drone pipelines.''}
\end{itemize}

\subsection{Slide 3: Engineering Philosophy}
\begin{itemize}
    \item \textbf{Focus:} Establish your standards.
    \item \textbf{Speaker Script:} 
    \textit{``My core engineering principles are: keeping pipelines decoupled, auditing dependencies, writing network-intercepted tests using MSW, and designing code that is ready for S3 and spatial indexes.''}
\end{itemize}

\subsection{Slide 4: STAGE 1: Sensor Data Ingestion (The Problem)}
\begin{itemize}
    \item \textbf{Focus:} How do we handle new image files?
    \item \textbf{Speaker Script:} 
    \textit{``Stage 1 begins when the airship returns from a flight. The raw image contains location and camera metadata. Our goal in Stage 1 is to read this telemetry locally on the client, render it on an interactive map, and handle fallbacks if the metadata is missing.''}
\end{itemize}

\subsection{Slide 5: Stage 1: Ingestion Stack}
\begin{itemize}
    \item \textbf{Focus:} Frontend tech choices.
    \item \textbf{Speaker Script:} 
    \textit{``The ingestion stack uses React and Vite. We use Leaflet for mapping, react-leaflet-cluster for UI optimization, exifr for client-side EXIF data parsing, and Axios to send payloads downstream.''}
\end{itemize}

\subsection{Slide 6: Stage 1: Component Layout}
\begin{itemize}
    \item \textbf{Focus:} UI structure.
    \item \textbf{Speaker Script:} 
    \textit{``The UI is split into five modules. MapPage manages the global state. MapView renders the Leaflet container, while UploadModal acts as the file gate and PhotoModal allows details inspection.''}
\end{itemize}

\subsection{Slide 7: Stage 1: Client EXIF Ingestion}
\begin{itemize}
    \item \textbf{Focus:} exifr code implementation.
    \item \textbf{Speaker Script:} 
    \textit{``When the user drops an image, we extract GPS headers using exifr locally in the browser before the upload begins. This is highly efficient because it removes image parsing overhead from the backend API.''}
\end{itemize}

\subsection{Slide 8: Stage 1: Advanced Camera Telemetry}
\begin{itemize}
    \item \textbf{Focus:} Telemetry data (altitude, yaw/pitch/roll).
    \item \textbf{Speaker Script:} 
    \textit{``For advanced inspections, we extract coordinate altitude and camera angles like yaw, pitch, and roll from XMP blocks. This allows us to calculate and project the exact camera view cone on the map.''}
\end{itemize}

\subsection{Slide 9: Stage 1: Manual Coordinate Overrides}
\begin{itemize}
    \item \textbf{Focus:} Fallback map click events.
    \item \textbf{Speaker Script:} 
    \textit{``If EXIF metadata is stripped, we display coordinate inputs and let the user click on an interactive mini-map. This captures coordinates instantly and ensures no photo is lost.''}
\end{itemize}

\subsection{Slide 10: Stage 1: Marker Clustering}
\begin{itemize}
    \item \textbf{Focus:} Map performance.
    \item \textbf{Speaker Script:} 
    \textit{``Rendering thousands of individual DOM markers crashes the browser. We group nearby coordinates using react-leaflet-cluster. Zooming in expands clusters dynamically, keeping map pan rates at 60 FPS.''}
\end{itemize}

\subsection{Slide 11: Stage 1: Visible Screen Bounding Boxes}
\begin{itemize}
    \item \textbf{Focus:} Viewport filtering (`moveend` listener).
    \item \textbf{Speaker Script:} 
    \textit{``To limit network payload sizes, the map captures Leaflet's moveend bounds. We query only coordinates inside the active screen window using debounced bounding box queries.''}
\end{itemize}

\subsection{Slide 12: STAGE 2: Backend API \& Ingestion Gateway}
\begin{itemize}
    \item \textbf{Focus:} Image transition from client to backend.
    \item \textbf{Speaker Script:} 
    \textit{``Now, the image crosses the network to Stage 2: The Gateway. Here, our Node/Express backend verifies the user session, parses the file using multer, normalizes the buffer, and commits metadata.''}
\end{itemize}

\subsection{Slide 13: Stage 2: Endpoint Routing Design}
\begin{itemize}
    \item \textbf{Focus:} REST endpoints structure and JWT guards.
    \item \textbf{Speaker Script:} 
    \textit{``Our REST routing is decoupled. Public routes manage sign-ins, while JWT auth middleware protects write actions like uploading, comment postings, and deletions.''}
\end{itemize}

\subsection{Slide 14: Stage 2: Relational Schema}
\begin{itemize}
    \item \textbf{Focus:} Schema design (Users, Photos, Comments).
    \item \textbf{Speaker Script:} 
    \textit{``I implemented a relational SQLite database schema. Deleting a photo triggers cascade rules to automatically clean up all child comment threads, maintaining referential integrity.''}
\end{itemize}

\subsection{Slide 15: Stage 2: Fast Lookup Indexing}
\begin{itemize}
    \item \textbf{Focus:} Database performance.
    \item \textbf{Speaker Script:} 
    \textit{``To support thousands of images, I created database indexes. I added a unique index on user emails, and spatial indexes on photo coordinates to speed up map queries.''}
\end{itemize}

\subsection{Slide 16: Stage 2: Image Normalization}
\begin{itemize}
    \item \textbf{Focus:} Sharp image orientation and compression.
    \item \textbf{Speaker Script:} 
    \textit{``Mobile uploads are often rotated. During upload, the backend routes the file to Sharp, calls rotate to fix orientation based on EXIF tags, and compresses the output to JPEG to save space.''}
\end{itemize}

\subsection{Slide 17: STAGE 3: Non-Blocking AI Processing}
\begin{itemize}
    \item \textbf{Focus:} Background jobs and Gemini Vision.
    \item \textbf{Speaker Script:} 
    \textit{``In Stage 3, the backend processes heavy tasks asynchronously. We save the database record, respond with a HTTP 201 immediately, and spawn a background AI worker to generate image descriptions.''}
\end{itemize}

\subsection{Slide 18: Stage 3: Vision AI Code Integration}
\begin{itemize}
    \item \textbf{Focus:} Gemini API call details.
    \item \textbf{Speaker Script:} 
    \textit{``Here is the code. We read the saved file buffer as a base64 string, invoke Gemini 1.5 Flash, request a descriptive caption, and write the output back to the database.''}
\end{itemize}

\subsection{Slide 19: Stage 3: AI Quota \& Retry Logic}
\begin{itemize}
    \item \textbf{Focus:} Exponential backoff logic.
    \item \textbf{Speaker Script:} 
    \textit{``If the AI API fails due to rate limits or drops, our code catches the error and executes an exponential backoff loop, retrying up to 3 times ($2\text{s} \to 4\text{s} \to 8\text{s}$) to ensure resilience.''}
\end{itemize}

\subsection{Slide 20: Stage 3: Manual AI Regeneration}
\begin{itemize}
    \item \textbf{Focus:} On-demand updates.
    \item \textbf{Speaker Script:} 
    \textit{``If the AI key was missing initially, users can click 'Regenerate' in the details view. This calls the backend update route and refreshes the UI state dynamically.''}
\end{itemize}

\subsection{Slide 21: STAGE 4: Security \& Quality Guardrails}
\begin{itemize}
    \item \textbf{Focus:} How we secured the pipeline.
    \item \textbf{Speaker Script:} 
    \textit{``In Stage 4, I audited and secured the pipeline. I replaced vulnerable random generators, wrote strict token lifetimes, and refactored dead code.''}
\end{itemize}

\subsection{Slide 22: Stage 4: Secure OTP Generation}
\begin{itemize}
    \item \textbf{Focus:} crypto.randomInt vs Math.random.
    \item \textbf{Speaker Script:} 
    \textit{``The original code used Math.random to generate verification codes. Since this is predictable, I refactored the backend to use Node's native crypto module to generate secure, high-entropy tokens.''}
\end{itemize}

\subsection{Slide 23: Stage 4: Expiration \& Revocation}
\begin{itemize}
    \item \textbf{Focus:} TTL and single-use map deletes.
    \item \textbf{Speaker Script:} 
    \textit{``Verification codes are stored in memory with a strict 10-minute expiry time-to-live. Codes are deleted immediately upon successful validation to block brute-force and replay attacks.''}
\end{itemize}

\subsection{Slide 24: Stage 4: Clean Code Refactoring}
\begin{itemize}
    \item \textbf{Focus:} Code cleanliness and file structures.
    \item \textbf{Speaker Script:} 
    \textit{``I cleaned up the codebase: removed dead state variables, removed informal messages from emails, and separated routes, middlewares, and helpers into clean folder structures.''}
\end{itemize}

\subsection{Slide 25: Stage 4: Environment Portability}
\begin{itemize}
    \item \textbf{Focus:} env variables.
    \item \textbf{Speaker Script:} 
    \textit{``I removed all hardcoded urls and keys. All routes and settings are managed through portable environment variables, allowing us to deploy the pipeline without changing code.''}
\end{itemize}

\subsection{Slide 26: STAGE 5: Pipeline Validation \& Testing}
\begin{itemize}
    \item \textbf{Focus:} How we prevent regressions in the pipeline.
    \item \textbf{Speaker Script:} 
    \textit{``In Stage 5, we validate our pipeline using automated tests. I set up unit tests, network mocking, package overrides, and CI/CD pipelines.''}
\end{itemize}

\subsection{Slide 27: Stage 5: DOM \& Leaflet Mocking}
\begin{itemize}
    \item \textbf{Focus:} setup.ts mocks.
    \item \textbf{Speaker Script:} 
    \textit{``JSDOM lacks browser APIs like localStorage and Leaflet map heights. In setup.ts, I mocked these layout APIs to ensure our tests compile and run reliably in console environments.''}
\end{itemize}

\subsection{Slide 28: Stage 5: MSW Network Mocking}
\begin{itemize}
    \item \textbf{Focus:} MSW handlers.
    \item \textbf{Speaker Script:} 
    \textit{``Instead of mocking Axios instances directly, I integrated Mock Service Worker to intercept requests at the browser network layer. This tests real network contracts without coupling tests to Axios.''}
\end{itemize}

\subsection{Slide 29: Stage 5: Dependency Security}
\begin{itemize}
    \item \textbf{Focus:} npm audit and overrides.
    \item \textbf{Speaker Script:} 
    \textit{``I ran npm audits, upgraded Vite to version 6 to fix a path traversal vulnerability, and configured overrides in package.json to pin secure esbuild versions.''}
\end{itemize}

\subsection{Slide 30: Stage 5: Docker \& CI/CD}
\begin{itemize}
    \item \textbf{Focus:} CI/CD pipelines.
    \item \textbf{Speaker Script:} 
    \textit{``I dockerized the project to lock Node runtimes. I also set up a GitHub Actions workflow that runs audits, formatting, and unit tests automatically on every Pull Request.''}
\end{itemize}

\subsection{Slide 31: STAGE 6: Scaling the Media Pipeline}
\begin{itemize}
    \item \textbf{Focus:} Production target for 10k+ images.
    \item \textbf{Speaker Script:} 
    \textit{``In the final stage, Stage 6, we scale the pipeline. We migrate our local SQLite database and file folders to production PostGIS and AWS S3 architectures.''}
\end{itemize}

\subsection{Slide 32: Stage 6: PostgreSQL \& PostGIS}
\begin{itemize}
    \item \textbf{Focus:} Why PostGIS.
    \item \textbf{Speaker Script:} 
    \textit{``SQLite is limited by file locking and lacks spatial indexes. In production, we migrate to PostgreSQL with PostGIS, storing coordinates in geography geom columns.''}
\end{itemize}

\subsection{Slide 33: Stage 6: GIST Spatial Indexing}
\begin{itemize}
    \item \textbf{Focus:} GIST grids.
    \item \textbf{Speaker Script:} 
    \textit{``Standard indexes cannot handle 2D map grids. We index geom columns using GIST, which groups coordinates into spatial bounding boxes, bypassing full-table scans.''}
\end{itemize}

\subsection{Slide 34: Stage 6: Viewport Envelope Queries}
\begin{itemize}
    \item \textbf{Focus:} ST_MakeEnvelope.
    \item \textbf{Speaker Script:} 
    \textit{``For viewport queries, we map map corners directly to PostGIS envelopes. The ST\_MakeEnvelope function and overlap operator check coordinate points inside boxes instantly.''}
\end{itemize}

\subsection{Slide 35: Stage 6: Proximity to Infrastructure}
\begin{itemize}
    \item \textbf{Focus:} ST_DWithin queries.
    \item \textbf{Speaker Script:} 
    \textit{``To automate inspection, we run spatial queries. Using ST\_DWithin, we query all photo coordinates taken within 50 meters of a specific pipeline route to link images to assets.''}
\end{itemize}

\subsection{Slide 36: Stage 6: Cloud Storage Architecture}
\begin{itemize}
    \item \textbf{Focus:} S3 Presigned URLs + CDN.
    \item \textbf{Speaker Script:} 
    \textit{``To handle thousands of uploads, clients bypass our API and upload files directly to S3 using presigned URLs. We then cache and distribute files through a CDN edge.''}
\end{itemize}

\subsection{Slide 37: Stage 6: Cookie Session Guards}
\begin{itemize}
    \item \textbf{Focus:} HttpOnly cookies & session security.
    \item \textbf{Speaker Script:} 
    \textit{``To protect user sessions from XSS attacks, we move JWTs out of LocalStorage into secure, HttpOnly cookies, implementing refresh token rotation on the backend.''}
\end{itemize}

\subsection{Slide 38: Retrospective \& Future Roadmap}
\begin{itemize}
    \item \textbf{Focus:} Limitations & what's next.
    \item \textbf{Speaker Script:} 
    \textit{``In retrospective, JSDOM map mocking was a major testing challenge. For the future, I plan to move map markers to MapLibre GL WebGL canvas and integrate YOLO to detect asset defects automatically on upload.''}
\end{itemize}

\subsection{Slide 39: Summary \& Open Q\&A}
\begin{itemize}
    \item \textbf{Focus:} Final wrap-up.
    \item \textbf{Speaker Script:} 
    \textit{``To summarize, GeoPhoto shows how we can build a decoupled, secure, and fully validated geospatial pipeline. Thank you for your time, and I am happy to open the floor to Q\&A.''}
\end{itemize}

\end{document}
